\documentclass[11pt,a4paper]{article}
\usepackage[T1]{fontenc}
\usepackage[utf8]{inputenc}
\usepackage[french]{babel}
\usepackage{lmodern}
\usepackage{amsmath,amssymb,mathtools}
\usepackage{geometry}
\geometry{top=22mm,bottom=22mm,left=23mm,right=23mm,headheight=15pt,headsep=6mm}
\usepackage{microtype,enumitem,booktabs,array,fancyhdr,lastpage,needspace}
\usepackage[hidelinks]{hyperref}
\setlength{\parindent}{0pt}
\setlength{\parskip}{6pt}
\setlist[enumerate]{leftmargin=*,label=\textbf{\alph*)},itemsep=8pt,topsep=4pt}
\pagestyle{fancy}\fancyhf{}
\renewcommand{\headrulewidth}{0.3pt}
\fancyfoot[C]{\small\thepage\,/\,\pageref{LastPage}}
\fancyfoot[L]{\scriptsize Version de travail -- 6 septembre 2026}
\setlength{\emergencystretch}{2em}
\hypersetup{pdftitle={MAT0339 -- Série C, examen 2 -- sujet},pdfauthor={Document de travail pédagogique}}
\fancyhead[L]{\small MAT0339 -- Série C -- Examen 2}
\fancyhead[R]{\small SUJET À VALIDER}
\begin{document}
{\LARGE\bfseries Série C -- Examen 2}\par
{\large Dénombrement, probabilités et algèbre linéaire}\par
\textbf{Épreuve étudiante -- version de travail}\par
Portée : chapitres 7 à 12. Total : 100 points. Durée proposée : 120 min.\par
\textbf{Nom :} \hrulefill\quad \textbf{Groupe :} \rule{22mm}{0.3pt}\par
\small Montrer les démarches. Conserver les valeurs exactes jusqu’à l’arrondi final; annoncer les unités. Convention provisoire : calculatrice scientifique non symbolique autorisée, sans notes. Les modalités doivent être confirmées par l’enseignant. Les pages suivantes offrent une page par question; une feuille supplémentaire peut être jointe.\normalsize\par
\vspace{3mm}\hrule\vspace{4mm}
{\large\bfseries Question 1 -- Compter les issues}\hfill\textbf{15 points}\par
\begin{enumerate}
\item \textbf{(5 points)} Un code comprend deux chiffres, suivis de deux lettres distinctes parmi les 26 lettres. Les chiffres peuvent se répéter et le zéro est permis. Combien de codes existent? Justifier la méthode.
\item \textbf{(5 points)} Combien de façons existe-t-il d’attribuer une présidence et un secrétariat à deux personnes distinctes parmi huit personnes? Justifier la méthode.
\item \textbf{(5 points)} Combien d’anagrammes distinctes peut-on former avec LETTRE? Toutes les lettres sont utilisées. Justifier la méthode.
\end{enumerate}
\vspace{3mm}\textit{Démarche et réponse :}\par\vfill

\newpage
{\large\bfseries Question 2 -- Tirages sans remise}\hfill\textbf{20 points}\par
\begin{enumerate}
\item \textbf{(8 points)} Une urne contient 3 boules rouges et 2 boules bleues. On tire uniformément deux boules successivement, sans remise. On note \(R_i\) et \(B_i\) la couleur au tirage \(i\). Construire un arbre pondéré et calculer les probabilités des quatre chemins.
\item \textbf{(4 points)} Calculer la probabilité d’obtenir deux boules de même couleur.
\item \textbf{(4 points)} Calculer \(P(R_1\mid B_2)\).
\item \textbf{(4 points)} Les événements \(R_1\) et \(R_2\) sont-ils indépendants? Justifier par une égalité ou une inégalité de probabilités.
\end{enumerate}
\vspace{3mm}\textit{Démarche et réponse :}\par\vfill

\newpage
{\large\bfseries Question 3 -- Essais répétés et espérance}\hfill\textbf{15 points}\par
\begin{enumerate}
\item \textbf{(4 points)} On effectue 6 essais indépendants, chacun de probabilité de succès \(p=0{,}5\). Pour \(X\), le nombre de succès, identifier le modèle et justifier les hypothèses.
\item \textbf{(4 points)} Calculer \(P(X=3)\).
\item \textbf{(3 points)} Calculer la probabilité d’au moins un succès.
\item \textbf{(4 points)} Pour un seul essai, le gain net est de 8~\$ en cas de succès et une perte de 2~\$ sinon. Calculer et interpréter l’espérance du gain net.
\end{enumerate}
\vspace{3mm}\textit{Démarche et réponse :}\par\vfill

\newpage
{\large\bfseries Question 4 -- Angle et projection}\hfill\textbf{15 points}\par
\begin{enumerate}
\item \textbf{(3 points)} Soient \(\mathbf u=(4,2),\quad\mathbf v=(0,3)\). Calculer les deux normes.
\item \textbf{(5 points)} Calculer le produit scalaire et l’angle entre les vecteurs, au centième de degré.
\item \textbf{(7 points)} Calculer \(\operatorname{proj}_{\mathbf v}\mathbf u\), puis vérifier que le vecteur résiduel est orthogonal à \(\mathbf v\).
\end{enumerate}
\vspace{3mm}\textit{Démarche et réponse :}\par\vfill

\newpage
{\large\bfseries Question 5 -- Composer et inverser des matrices}\hfill\textbf{20 points}\par
\begin{enumerate}
\item \textbf{(6 points)} Soient \[A=\begin{pmatrix}3 & 1\\2 & 1\end{pmatrix},\qquad B=\begin{pmatrix}1 & 0\\1 & 1\end{pmatrix}.\] Calculer \(AB\) et \(BA\). Comparer les résultats.
\item \textbf{(6 points)} Calculer \(\det A\), puis \(A^{-1}\). Vérifier le produit \(AA^{-1}\).
\item \textbf{(8 points)} Résoudre \(A\mathbf x=\begin{pmatrix}11\\8\end{pmatrix}\). Vérifier dans les deux équations initiales.
\end{enumerate}
\vspace{3mm}\textit{Démarche et réponse :}\par\vfill

\newpage
{\large\bfseries Question 6 -- Intersection et classification d’un système}\hfill\textbf{15 points}\par
\begin{enumerate}
\item \textbf{(8 points)} La droite \((x,y,z)=(2,-1,0)+t(1,1,2)\), \(t\in\mathbb R\), rencontre-t-elle le plan \(x-y+z=7\)? Déterminer leur intersection et justifier son unicité.
\item \textbf{(7 points)} En fonction du réel \(k\), classer les solutions du système \[\begin{cases}x+y+z=5,\\x-y+z=1,\\2x+2z=k.\end{cases}\] Donner une paramétrisation lorsqu’il est compatible.
\end{enumerate}
\vspace{3mm}\textit{Démarche et réponse :}\par\vfill
\end{document}
